% Part: computability
% Chapter: recursive-functions
% Section: primitive-recursion

\documentclass[../../../include/open-logic-section]{subfiles}

\begin{document}

\olfileid[hi]{cmp}{rec}{pre}
\olsection{आदिम पुनरावर्तन}

प्राकृतिक संख्याओं की एक विशेषता यह है कि~$0$ से आरंभ करके उत्तराधिकारी
संक्रिया~$+1$ को परिमित बार लागू करने पर प्रत्येक प्राकृतिक संख्या तक पहुँचा जा सकता
है---हर प्राकृतिक संख्या या तो~$0$ है, या \dots{}~$0$ के उत्तराधिकारी
का उत्तराधिकारी है। इस तथ्य का उपयोग करने वाले किसी फलन
~$h\colon\Nat \to \Nat$ को निर्दिष्ट करने का एक ढंग यह है: (a)~यह
निर्दिष्ट करें कि $h$ का तर्क~$0$ पर मान क्या है, और (b)~यह भी निर्दिष्ट
करें कि $h(x)$ का मान दिए होने पर $h(x+1)$ का मान कैसे संगणित किया जाए।
क्योंकि (a) हमें सीधे बताता है कि $h(0)$ क्या है, इसलिए $h$,~$0$ पर
परिभाषित है। अब $x=0$ के लिए (b) में दिया निर्देश लेकर हम~$h(1) = h(0+1)$ को
$h(0)$ से संगणित कर सकते हैं। $x=1$ के लिए उन्हीं निर्देशों से हम
~$h(2) = h(1+1)$ को $h(1)$ से संगणित करते हैं, और इसी प्रकार आगे बढ़ते हैं।
हर प्राकृतिक संख्या~$x$ के लिए अंततः हम उस चरण तक पहुँचेंगे जहाँ $h(x)$
से $h(x+1)$ परिभाषित करते हैं; अतः $h(x)$ प्रत्येक $x \in \Nat$ के
लिए परिभाषित है।

मसलन, मान लें कि हम निम्न दो समीकरणों से $h\colon \Nat \to \Nat$
निर्दिष्ट करते हैं:
\begin{align*}
h(0) & =  1\\
h(x+1) & =  2 \cdot h(x)
\end{align*}
यदि हम गुणा करना पहले से जानते हैं, तो ये समीकरण हमें ऊपर के (a) और~(b)
के लिए आवश्यक सूचना देते हैं। दूसरा समीकरण क्रमशः लागू करने पर मिलता है:
\begin{align*}
  h(1) & = 2\cdot h(0) = 2,\\
  h(2) & = 2\cdot h(1) = 2\cdot 2,\\
  h(3) & = 2 \cdot h(2) = 2\cdot 2 \cdot 2,\\
  & \vdots
\end{align*}
हम देखते हैं कि हमारे द्वारा निर्दिष्ट फलन~$h$, $h(x) = 2^x$ है।

प्राकृतिक संख्याओं की यह विशेषता सुनिश्चित करती है कि इन दोनों कसौटियों
को पूरा करने वाला केवल एक फलन~$h$ है। ऐसे समीकरणों के युग्म को फलन~$h$
की \emph{आदिम पुनरावर्तन द्वारा परिभाषा} कहते हैं। ऐसा इसलिए कहते हैं कि
हम~$h$ को ``पुनरावर्ती रूप से'' परिभाषित करते हैं, अर्थात् परिभाषा में,
विशेषकर दूसरे समीकरण के दाहिने पक्ष में, स्वयं $h$ आता है। यह ``आदिम''
है क्योंकि $h(x+1)$ को परिभाषित करने में हम केवल मान~$h(x)$, अर्थात् ठीक
पूर्ववर्ती मान, का उपयोग करते हैं।~$\Nat$ पर किसी फलन को पुनरावर्ती रूप
से परिभाषित करने का यह सबसे सरल ढंग है।

जोड़ और गुणा जैसे और भी मूलभूत फलनों को हम आदिम पुनरावर्तन से परिभाषित
कर सकते हैं। किंतु इन स्थितियों में संबंधित फलन $2$-स्थानी हैं। हम एक
तर्क-स्थान को नियत करते हैं और दूसरे का उपयोग पुनरावर्तन के लिए करते हैं।
मसलन, $\Add(x, y)$ को परिभाषित करने के लिए हम~$x$ को नियत करके पहले
$y=0$ पर मान और फिर $y+1$ पर मान को~$y$ के पदों में परिभाषित कर सकते हैं।
चूँकि $x$ नियत है, वह परिभाषक समीकरणों के बाएँ और दाएँ दोनों पक्षों में
आएगा।
\begin{align*}
\Add(x,0) & =  x\\
\Add(x,y+1) & =  \Add(x,y)+1
\end{align*}
ये समीकरण $\Add$ का मान सभी $x$ \emph{और}~$y$ के लिए निर्दिष्ट करते
हैं। मसलन, $\Add(2,3)$ ज्ञात करने के लिए हम $x = 2$ पर परिभाषक समीकरण
लागू करते हैं: पहले समीकरण से $\Add(2,0) = 2$ और फिर दूसरे से क्रमशः
$\Add(2,1) = 2 + 1 = 3$, $\Add(2, 2) = 3 + 1 = 4$, $\Add(2,
3) = 4 + 1 = 5$ मिलता है।

$\Add$ की परिभाषा में हमने दूसरे समीकरण के दाहिने पक्ष पर $+$ का उपयोग
किया, किंतु केवल~$1$ जोड़ने के लिए। दूसरे शब्दों में, हमने उत्तराधिकारी
फलन $\Succ(z) = z+1$ का उपयोग किया और उसे पूर्ववर्ती मान
$\Add(x,y)$ पर लागू करके $\Add(x,y+1)$ परिभाषित किया। अतः हम पुनरावर्ती
परिभाषा को ऐसे एक फलन के पदों में दी हुई समझ सकते हैं जिसे हम पूर्ववर्ती
मान पर लागू करते हैं। फिर भी फलन को केवल पूर्ववर्ती मान पर नहीं, बल्कि
$x$ और~$y$ पर भी निर्भर होने देना हानिकर नहीं---और कभी-कभी आवश्यक---है।
विचार करें:
\begin{align*}
  \Mult(x,0) & =  0 \\
  \Mult(x,y+1) & =  \Add(\Mult(x,y),x)
\end{align*}
यह फलन $\Mult$ की आदिम पुनरावर्ती परिभाषा है, जिसमें फलन $\Add$ को
पूर्ववर्ती मान $\Mult(x,y)$ और प्रथम तर्क~$x$, दोनों पर लागू किया जाता है।
यह फलन~$\Mult(x,y)$ को भी सभी तर्कों $x$ और~$y$ के लिए परिभाषित करता है।
मसलन, $\Mult(2,3)$ को $\Mult(2,0)$, $\Mult(2,1)$, $\Mult(2,2)$ और
~$\Mult(2,3)$ क्रमशः संगणित करके निर्धारित किया जाता है:
\begin{align*}
  \Mult(2,0) & = 0\\
  \Mult(2,1) & = \Mult(2,0+1) =
  \Add(\Mult(2,0), 2) = \Add(0, 2) = 2\\
  \Mult(2,2) & = \Mult(2,1+1) =
  \Add(\Mult(2,1), 2) = \Add(2, 2) = 4\\
  \Mult(2,3) & = \Mult(2,2+1) =
  \Add(\Mult(2,2), 2) = \Add(4, 2) = 6
\end{align*}

तब सामान्य प्रतिरूप यह है: किसी फलन~$h(x_0, \dots, x_{k-1}, y)$ की
आदिम पुनरावर्ती परिभाषा देने के लिए हम दो समीकरण देते हैं। पहला,
$h(x_0, \dots, x_{k-1}, 0)$ का मान~$h$ का संदर्भ लिए बिना परिभाषित करता है।
दूसरा, $h(x_0,
\dots, x_{k-1}, y+1)$ का मान
$h(x_0, \dots, x_{k-1}, y)$, अन्य तर्कों $x_0$,
\dots,~$x_{k-1}$ और~$y$ के पदों में परिभाषित करता है। उस दूसरे समीकरण
में~$h$ के केवल ठीक पूर्ववर्ती मान
का उपयोग किया जा सकता है। यदि हम इन दोनों समीकरणों के दाहिने पक्षों से
दी गई संक्रियाओं को स्वयं फलन $f$ और~$g$ मानें, तो आदिम पुनरावर्तन द्वारा
नया फलन~$h$ परिभाषित करने का सामान्य प्रतिरूप यह है:
\begin{align*}
  h(x_0, \dots, x_{k-1}, 0) & = f(x_0, \dots, x_{k-1})\\
  h(x_0, \dots, x_{k-1}, y+1) & =
  g(x_0, \dots, x_{k-1}, y, h(x_0, \dots, x_{k-1}, y))
\end{align*}  
$\Add$ की स्थिति में $k=1$ और $f(x_0) = x_0$ (तत्समक फलन) है, तथा
$g(x_0, y, z) = z + 1$ (वह $3$-स्थानी फलन जो अपने तीसरे तर्क का
उत्तराधिकारी लौटाता है) है:
\begin{align*}
  \Add(x_0, 0) & = f(x_0) = x_0\\
  \Add(x_0, y+1) & = g(x_0, y, \Add(x_0, y)) =
  \Succ(\Add(x_0, y))
\end{align*}
$\Mult$ की स्थिति में $f(x_0) = 0$ (सदा~$0$ लौटाने वाला अचर फलन) और
$g(x_0, y, z) = \Add(z,x_0)$ (अपने अंतिम और प्रथम तर्कों का योग लौटाने
वाला $3$-स्थानी फलन) है:
\begin{align*}
  \Mult(x_0,0) & =  f(x_0) = 0 \\
  \Mult(x_0,y+1) & =  g(x_0, y, \Mult(x_0,y)) =
  \Add(\Mult(x_0,y), x_0)
\end{align*}

\end{document}
